%=====================================================================
%%  _____      _                 _            _   _
%% |_   _|    | |               | |          | | (_)
%%   | | _ __ | |_ _ __ ___   __| |_   _  ___| |_ _  ___  _ __
%%   | || '_ \| __| '__/ _ \ / _` | | | |/ __| __| |/ _ \| '_ \
%%  _| || | | | |_| | | (_) | (_| | |_| | (__| |_| | (_) | | | |
%%  \___/_| |_|\__|_|  \___/ \__,_|\__,_|\___|\__|_|\___/|_| |_|
%=====================================================================
 \chapter{Introduction}
\label{cha:introduction}
\fancyhead[LO]{Introduction}
\fancyhead[ER]{Introduction}
%=====================================================================
%=====================================================================

The simultaneous flow of materials with different components or
different states of matter such as gas, liquid, solid, and plasma is
generally called multiphase flow.  The simplest case of a multiphase
flow is the two-phase flow consisting of: (1) two states of the same
component (e.g., water and steam), (2) two components of the same
state
(e.g., oil and water), or (3) two components of two different states
(e.g., solid particles in gas). Computational methods have become a
key tool for investigating the fundamentals of such flows allowing to
explore the underlying physics in great detail and on different
scales. Consequently, numerical simulations of multiphase flows are
nowadays extensively employed to tackle engineering problems in
several industrial branches.

The scope of the present thesis is the numerical modeling and
simulation of turbulent fluid flows involving cohesive solid
particles. In particular, the focus is directed towards the breakage
of agglomerates, which plays an important role in dictating the
particle size distribution in many particle-laden flows.  The
effective decrease of the particle size by the disintegration of
agglomerates is in favor of obtaining desirable material properties
and operational conditions, e.g., delivering medical aerosols with
uniform dosing. On the contrary, breakage can be problematic as the
reduction of the particle size decreases the settling velocity of
particles in cyclone separators hindering the separation process.

In this chapter some examples pointing out the significance of
particle-laden flows in general will be discussed first, followed by
an overview on the most important modeling approaches available in the
literature (Section~\ref{sec:modeling_approaches}). In
Section~\ref{sec:aims_objectives} the objectives of this thesis are
elaborated and clearly defined. Lastly, the organization of the thesis
is outlined in Section~\ref{sec:Outline}.

% The motivation
%for studying particle-laden flows is given in
%Section~\ref{sec:motiv},
%followed by an
%overview on
%the available numerical modeling approaches
%in
%Section~\ref{sec:modeling_approaches}. In
%Section~\ref{sec:aims_objectives} the aims of this thesis are
%elaborated and the objectives are clearly defined. Lastly, the
%organization of the topics in this thesis is outlined in
%Section~\ref{sec:Outline}.

%The interactions
%within the single phase and between the two phases give rise to
%numerous complex yet interesting phenomena, which have been the
%the subject of numerous experimental and numerical studies
%throughout the last decades
%=====================================================================
%=====================================================================
\section{Motivation for Considering Particle-Laden Flows}
\label{sec:motiv}
\fancyhead[OL]{Motivation for Considering Particle-Laden Flows}
%=====================================================================
%=====================================================================

Particle-laden flows are abundant in nature. Typical examples of
environmental processes include the formation of dust and sand storms,
the transportation of sediment in rivers and streams, and the
formation of ash clouds by volcanic eruption plumes. The satellite
image shown in Fig.~\ref{subfig:sandstorm} depicts clouds of sand and
dust originated from the Sahara (north Africa) and transported to the
European continent across the Mediterranean Sea with a span of about
10$^3$ km.  Such dust storms deteriorate the air quality, worsen the
public health, and affect the frequency of North Atlantic hurricanes
\citep{clifford2019}. Considering that the size of sand particles
varies between microns and millimeters, the image portrays a special
characteristic of particle-laden flows, which is the potential vast
range of spatial and temporal scales involved in the physics of a
single problem.  Another related example from nature is the Zambezi
river delta depicted in Fig.~\ref{subfig:zambeszi}, which flows about
230 km across six African countries to eventually pour its water and
transported sediment into the Indian Ocean. This discharge of
sediments preserves the quality and softness of the water inside the
Delta, leading to fertile soil and suitable water for farming and
fishing, respectively.  Unfortunately, the Zambeszi river nowadays is
not capable of transporting the sediment out of the continent like it
used to do due to the presence of hydropower dams upstream, which in
turn reduce the flow of the river causing the sediment to accumulate
\citep{disskunz2011effect}.

%---------------------------------------------------------------------
\begin{figure}
\label{fig:particle_laden_exam}
\centering
%
\subfloat[
	\label{subfig:sandstorm}]{
	\includegraphics[draft=false,width=0.485\linewidth,trim=
 0.0cm 0cm 0cm 0.0cm,clip]{figure/intro/sdusteurope_ann.eps}}%
\hfill
%
\subfloat[
	\label{subfig:zambeszi}]{
	\includegraphics[draft=false,width=0.485\linewidth,trim=
 0.0cm 0.0cm 0.0cm 0.0cm,clip]{figure/intro/zambezi_ann.eps}}%
\\ \vspace{0.5cm}
%
\hspace{2cm}
\subfloat[
	\label{subfig:dpi}]{
	\includegraphics[draft=false,width=0.26\linewidth]{figure/intro/easyinhaler2.eps}}%
\hfill
%
\subfloat[
	\label{subfig:pollutant_fac}]{
	\includegraphics[draft=false,width=0.49\linewidth]{figure/intro/Smoke2.jpg}}
%
\caption{Examples on particle-laden flows: (a) A dust storm traveling
  from north Africa to Italy spanning about 10$^3$ km, (b) The
  Zambezi river in Africa with a total length of 2574 km
  pouring water and sediment into the Indian Ocean (Images: NASA,
  Visible Earth), (c) A dry powder inhaler
  \textit{Easyhaler\textsuperscript{\sffamily\textregistered}}
  \citep[after][]{easyhaler}, (d) Exhaust fumes from smoking
  chimneys (Image:
  University of Freiburg).}
\end{figure}
%\texttt{http://visibleearth.nasa.gov}
%---------------------------------------------------------------------

Besides nature, particle-laden flows are present in various technical
applications ranging from cyclone separators, particle-based solar
receivers, pharmaceutical inhalers (see Fig.~\ref{subfig:dpi}) to the
emission of pollutants in different industrial branches (see
Fig.~\ref{subfig:pollutant_fac}). In many of these applications,
particles are found to be naturally cohesive and the particle size is
normally measured in microns. In this specific size range, the effect
of cohesion is found to be significant and thus the presence of
particles in the form of clusters (i.e., agglomerates) is very likely.
However, agglomerates in turn are prone to breakage (i.e.,
fragmentation) due to intense interactions with the turbulent flow or
strong collisions with the surrounding objects. These circumstances
are fulfilled in many applications since turbulent, wall-bounded flows
(e.g., chimney, mouth piece of an inhaler) are dominant in
practice. Hence, the size evolution of the particle phase in such
systems is governed by a competition between counteracting physical
phenomena.

%---------------------------------------------------------------------
\begin{figure}
\centering
\includegraphics[draft=false,width=0.6\textwidth]{figure/intro/new_lung.eps}%
\caption{Inhalation of particulates
in the human respiratory tract  \citep[after][]{leimao2011}.}
\label{fig:lung}
\end{figure}
%\texttt{http://visibleearth.nasa.gov}
%---------------------------------------------------------------------

In fact, the particle size is a crucial feature for the aforementioned
applications and others. For instance, the particle size plays a
major role for optimizing the efficacy of inhaled medications or
describing the harmfulness of particulate pollutants.
Fig.~\ref{fig:lung} schematically depicts particles deposited in the
human respiratory system. As illustrated, particles are able to reach
different regions of the respiratory tract depending on their size
among other factors \citep{darquenne2012aerosol}. In general, large
particles \mbox{(> 6 $\mu$m)} are unable to exactly follow the paths
of the inhaled air when a bend is met due to their high
inertia. Consequently, large particles impact and deposit in the upper
airways, e.g., in the nose or throat and thus drugs cannot be
delivered
to
the lung tissues. In order to effectively treat the central and narrow
airways, particles must be in the size range of 2 to 6 $\mu$m. Even
smaller
particles tend to penetrate deep into the
alveolar region where the exchange of oxygen and carbon dioxide takes
place or even into the bloodstream posing great health problems in
case of pollutants.


In summary, particle-laden flows are of paramount importance to
diverse fields of science and engineering. The examples discussed
above indicate that the profound knowledge of such flows can improve
the human well-being through the optimization of medical and other
consumable products and contribute  to the alleviation of serious
environmental problems. In many cases, the design and the dynamics of
particle-laden flows strongly depends on the particle size and its
evolution within the system.  Therefore, it is inevitable to take
phenomena such as agglomeration and breakup in predictions into
account to obtain
useful and physically meaningful forecasts. %Experimental research and
%numerical
%simulations are

%    \AKp{importance for every day life and industry and complex
%    nature due to wide range of scales involved. experiment and
%    numerical tool are needed. the focus is on the breakage and
%    agglomeration phenomenon}
%the breakage of solid particle agglomerates dispersed in turbulent
%flows.
%Furthermore, the light is shed on the
%interplay between breakage and agglomeration in flow scenarios
%relevant in practice.
%=====================================================================
%=====================================================================
\section{Modeling Approaches}
\label{sec:modeling_approaches}
\fancyhead[OL]{Modeling Approaches}
%=====================================================================
%=====================================================================

A major challenge confronted in the computational modeling of
turbulent particle-laden flows is the extensive range of length and
time scales to be dealt with \citep{balachandar2009scaling}. This is
because turbulence, which is inherently a multiscale phenomenon, needs
to be described while additionally accounting for fluid-particle and
inter-particle interactions occurring on various scales. Another
challenge is posed by the inertial motion of relatively large and
heavy particles, which can strongly deviate from the streamlines of
the carrier fluid \citep{disscapecelatro2014}. The complexity of such
flows inspired the development of different modeling approaches in the
literature, which compete through the level of captured details and
the computational expenses.  Fig.~\ref{fig:methods} schematically
illustrates the main features of the different approaches, which will
be briefly explained in the following.

 %---------------------------------------------------------------------
 \begin{figure}
 	\centering
 	\includegraphics[draft=false,width=1\textwidth]{figure/intro/methods.eps}%
 	\caption{Schematic representation of the common  modeling
 	approaches for
 	particle-laden flows
 	\citep[after][]{capecelatro2021modeling}.}
 	\label{fig:methods}
 \end{figure}
 %---------------------------------------------------------------------

 The optimum modeling approach for a certain problem is traditionally
 postulated based on a regime map distinguishing particle-laden flow
 problems according to two features
 \citep{elghobashi1991,elghobashi1994,balachandar2009scaling,balachandar2010}.
 The first is the volume fraction of the solid phase, expressing how
 dilute or dense a system is. The second feature is the particle size
 or the particle response time relative to the corresponding smallest
 turbulent scales. However, efforts are being continually undertaken
 in the literature to extend the range of applicability of different
 modeling approaches
 \citep[e.g.,][]{balachandar2020towards,masi2014development,akiki2017pairwisepar,evrard2020euler}.
 Hence, the limits defining the validity range of these methods are
 relaxed leading to an increased flexibility in the decision making
 process.  In the following, the attention is restricted to describing
 the most important features of the numerical modeling approaches used
 in recent research. \newline

The first approach is the Euler--Euler (two-fluid model) formulation
treating the continuous and the disperse phases as interpenetrating
continua.  Consequently, conservation equations need to be solved for
both phases. These are most frequently expressed based on the
Reynolds-Averaged Navier-Stokes equations (RANS) delivering
statistical descriptions of both phases
\citep[e.g.,][]{elghobashi1983two,zhou2010advances}. Nevertheless,
different studies have been recently devoted to develop instantaneous
Euler--Euler methods in the context of large-eddy simulations (LES)
\citep[e.g.,][]{moreau2010development,ma2015euler} and direct
numerical simulations (DNS) \citep[e.g.,][]{masi2014development}. In
the Euler--Euler approach the interactions between particles are not
realized in terms of discrete collision events. Instead, the momentum
transfer within the solid phase is modeled through the stress tensor
appearing in the momentum conservation equation relying on an analogy
with the fluid phase \citep{van2001comparative}.  In addition, the two
phases are coupled by an exchange of momentum and energy. For this
purpose, source terms described by constitutive models are introduced
into the corresponding conservation equations
\citep{balachandar2009scaling}.

The Euler--Euler approach has been widely applied in simulations
involving a huge number of particles since the computational load is
related to the solution of the governing equations and does not scale
with the number of particles \citep{sommerfeld2008}. However, if
particles possessing different sizes are present (i.e., a polydisperse
system), each size class is ideally treated as a different continuum
(multi-fluid model). For this purpose, a set of conservation equations
needs to be solved for each size class requiring considerably higher
computational efforts \citep{kuerten2016}.  In addition, predicting
the development of the particle size due to agglomeration and breakage
requires coupling the Euler--Euler approach with CPU-time intensive
population balance models
\citep{ramkrishna2000population,marchisio2003quadrature}. The solution
of the population balance equations depends on collision,
agglomeration, and breakage rate functions, which are often termed as
kernels. These kernels are mostly derived from experiments relying on
certain assumptions regarding the flow and particle properties.  That
renders the Euler--Euler solution in real applications heavily
dependent on modeling issues since the discrete nature of the
particles is ignored. \newline

The other extreme is the fully resolved approach where the flow scales
introduced by the particle beside all the scales of the ambient
turbulence are completely resolved. Various methods can be used to
capture the unsteady boundary layer around the particles and the
resulting wakes. In the context of Navier-Stokes solvers, direct
numerical simulations based on body-fitted spherical grids or on the
immersed boundary method are common techniques
\citep{burton2005fully,uhlmann2008,vreman2016particle,vowinckel2019consolidation,vowinckel2019settling,uhlmann2005}.
As an alternative, fully resolved simulations can be conducted
applying the lattice-Boltzmann method
\citep{cui2018application,jebakumar2019fully,dietzel2016application}.
Obviously, the fully resolved approach offers the most detailed
description of the physics behind the fluid-particle interactions.
Consequently, the motion of particles follows directly from Newton's
second law using the resolved forces and torques imposed by the
ambient fluid.  However, despite the ever-growing available
computational power, such methods are still restricted to dilute
systems involving up to $\mathcal{O} (10^4)$ particles
\citep[e.g.,][]{de2019assessment,picano2015turbulent,jain2021impact,uhlmann2017clustering,vowinckel2014fluid,kempe2014relevance,kempe2012collision}.
Even for such relatively small systems the problem becomes quickly
intractable if agglomerates are present and the inter-particle
interactions should be resolved in space and time. Hence, at the
present time fully resolved simulations are not an affordable option
to comprehensively investigate real applications.  Nevertheless, the
insightful knowledge offered by the high resolution is frequently
exploited to derive and improve closure models for the less detailed
approaches. This is typically achieved by conducting resolved
simulations in controlled configurations while focusing on a specific
phenomenon, e.g., the drag force and the lift force
\citep[e.g.,][]{beetstra2007drag,bogner2015drag,
  yali2015new,akiki2020shear,kurose99,bagchi2002shear,holzer2009lattice,dietzel2016application}.
Recently, the utilization of detailed results in improving efficient
but less conclusive models has been greatly empowered by data-driven
approaches relying on modern machine learning techniques
\citep{montans2019data,duraisamy2019turbulence,beck2019deep,balachandar2020towards}.
In particular, artificial neural networks have shown great potential
to efficiently reproduce knowledge obtained under realistic operating
conditions. That is a clear advantage over traditional models, which
mostly refer to simplified scenarios. \newline

The third approach is the Euler--Lagrange point-particle method, which
finds a middle ground between the two extremes (Euler--Euler and fully
resolved particle). In this framework the carrier fluid is treated as
usual in the Eulerian frame of references. In addition, a Lagrangian
tracking based on Newton's second law is applied to describe the
motion of particles. However, in contrast to the fully resolved
approach the particles are approximated as point-masses. Thus, the
fluid-particle interface vanishes and the motion of particles is
assumed to be governed by the local properties of the undisturbed flow
at the particle center \citep{maxey83}. For this purpose, a variety of
closure models has been proposed in the literature to describe the
forces and torques exerted by the fluid in the equation of motion of
the particle. If the volume fraction of the disperse phase is
sufficiently large, it becomes necessary to account for the feedback
effect of the particles on the fluid (i.e., two-way coupling). This is
typically done by adding a source term to the momentum conservation
equation of the fluid.

Due to the circumstance that the small flow scales induced by the
presence of particles are ignored, the ambient turbulence does not
necessarily need to be fully resolved by time-consuming DNS. Depending
on the required accuracy and the available computational resources
alternatively, LES or RANS can be employed to predict the turbulent
flow field \citep{balachandar2009scaling,kuerten2016}. The latter can
offer a drastic reduction of the computational costs especially for
statistically steady turbulent flow problems. This is attributed to
the fact that the RANS solution solely provides an averaged
description of turbulence allowing to "freeze" the mean flow field
once the solution is developed. As a result, the tracking of particles
follows as a post-processing step \citep{kuerten2016}.  However, the
challenge here is that particles in reality interact with the
instantaneous flow field and not the averaged one.  Therefore, the
velocity fluctuations need to be mimicked by additional models of
stochastic nature \citep{sommerfeld2008}. For these reasons, LES is
gaining a considerable attention since it offers reasonably accurate
instantaneous solutions by resolving the most important flow
structures at moderate costs.  However, often the effects of the
unresolved subgrid-scale structures on the flow field and the
particles are not negligible and thus need to be adequately accounted
for.

Furthermore, the Euler--Lagrange method offers the possibility to
describe the inter-particle interactions (i.e., four-way coupling) at
different levels of detail. Owing to the persistence of the discrete
nature of particles in this method, the collision events can be
individually detected either by stochastic \citep{sommerfeld2001} or
deterministic \citep{bird76,breuer2012b} algorithms. For particle-wall
collisions the detection is even more straightforward.

 Collisions lay the groundwork for important phenomena such as
 particle agglomeration, particle-wall adhesion, and agglomerate
 breakage. This implies that the prediction of these phenomena
 substantially depends on the method by which the collisions are
 treated. For this purpose, two methods are commonly applied. The
 first is the soft-sphere approach, which is also known as the
 discrete element method (DEM) \citep{cundall1979}. This method
 describes collisions as gradually evolving processes which are
 resolved in space and time. The interactions between a certain
 particle and multiple others can be simultaneously taken into
 account.  Consequently, particles form agglomerates and deposit at
 walls according to the balance of forces between the interacting
 objects. Hence, agglomeration and wall-adhesion are implicitly taken
 into account without additional modeling.  Following the same
 concept, the breakup of agglomerates and the re-suspension of
 deposited particles naturally occurs. However, this direct prediction
 is associated with high computational burdens since apart from the
 costly full resolution of the collisions, the tracking of contacts is
 not trivial. A cost-efficient alternative for the soft-sphere
 treatment is the hard-sphere model which treats collisions as instant
 events instead of resolving them in time \citep{crowe98}. In the
 original formulation of this model solely the post-collision dynamics
 are predicted by a momentum balance employing the pre-collision
 conditions of the colliding objects. Agglomeration and wall adhesion
 can also be predicted either by extending the underlying momentum
 balance \citep{kosinski2009,kosinski2010,breuer2015a,almohammed2016b}
 or by applying energy-based models
 \citep{disshiller1981,dissjuergens2012,almohammed2016}. A drawback of
 the hard-sphere model is that it solely handles binary collisions.

  Agglomerates may be geometrically represented by their real,
  fracture-like structures. In this case, a physically meaningful
  description of the dynamics of the agglomerates requires a
  continuous determination of the fluid-particle and particle-particle
  interactions for each primary particle.  In another approach, the
  structure of agglomerates is replaced by a single particle
  possessing an effective diameter, which allows to track
  agglomerates as single entities.  The discussion on both concepts
  including the possible methods to treat the fluid-particle and
  particle-particle interactions and their implications for the
  prediction of the agglomerate breakage is extended in
  Chapter~\ref{cha:literature_review}. \newline

In summary, a variety of computational methods exist to describe
particle-laden flows. The Euler--Lagrange approach combined with the
point-particle approximations offers important advantages over the
other methods. On the one hand, compared to Euler--Euler models, it
allows a greater flexibility in the level of captured details for both
phases and their coupling.  On the other hand, it offers a significant
reduction of the computational costs (multiple orders of magnitude)
over fully-resolved models. Consequently, simulating flows comprising
millions of particles while accounting for diverse flow and particle
phenomena is possible.  However, efforts are still required in order
to further narrow the accuracy gap between the point-particle models
and the fully resolved solutions.

%=====================================================================
%=====================================================================
\section{Objectives of the Thesis}
\label{sec:aims_objectives}
\fancyhead[OL]{Objectives of the Thesis}
%=====================================================================
%=====================================================================

The main goal of this thesis is to advance the current state of the
art in computational modeling of particle-laden flows in an important
yet less-established area.  More specifically, the scope is to develop
physical models for numerical simulations describing the breakage of
particle agglomerates in turbulent flows. The models are intended for
(but not restricted to) the in-house CFD code \lesocc \ (Large-Eddy
Simulation On Curvilinear
Coordinates\nomenclature[A]{\lesocc}{Large-Eddy Simulation On
  Curvilinear Coordinates})
\citep{breuer1998,breuer1998b,breuer2000,habil2002}, which provides an
efficient multiphase flow simulation methodology relying on large-eddy
simulation for the continuous phase and a Lagrangian point-particle
tracking for the disperse phase.  The aim is to propose an advanced
simulation strategy capable of delivering reliable predictions of
particle-laden flows in practically relevant setups. This can only be
achieved by reasonably taking various important phenomena at
manageable computational costs into account.

In more practical terms, the objectives of the present thesis are to:

\begin{enumerate}[label=(\arabic*),align=left, leftmargin=*]
	\item Enhance the modeling of the structural features of
          agglomerates, namely, the porosity and the coordination
          number (i.e., average number of contacts per particle).

	\item Develop a model for the breakage of agglomerates by
          fluid-induced stresses.

	\item Develop a model for the breakage of agglomerates by wall
          impacts.

     \item Apply the entire simulation methodology including the newly
       developed models to investigate the breakage of agglomerates in
       wall-bounded turbulent flows under various operating
       conditions.
\end{enumerate}

For the first objective, an adequate representation of the porosity
and the coordination number is sought. Since agglomerates are porous
by nature, the porosity is an essential characteristic for determining
the corresponding density and diameter of the spherical particles used
to replace the structure of agglomerates in the present
methodology. In addition, both the porosity and the coordination
number are important parameters for determining the resistance of
agglomerates against breakage, i.e., its strength.

Simplifying the complex structure of agglomerates by single spheres
poses certain challenges on the prediction of breakage, which need to
be tackled. More specifically, the full characterization of a breakage
event in the second and third objectives requires addressing three
issues, hereafter called the breakage modeling tasks:

\begin{enumerate}[label=(\roman*),align=left, leftmargin=*]
 	\item A criterion for the onset of breakage.
 	\item The number of resulting fragments and their size
          distribution.
 	\item The velocity vectors of the fragments arising due to
          breakage.
 \end{enumerate}

The attention in this thesis is restricted to spherical, monodisperse
(i.e., equally-sized) primary particles in the micrometer size range.
In addition, the particles are assumed to be dry and electrostatically
neutral, and the cohesion between particles can solely be attributed
to the van-der-Waals force. All other specific assumptions will be
mentioned in the corresponding discussion.

To realize the fourth objective, test cases inspired by experimental
investigations reported in the literature are taken into account. The
aim is to assess the influence of a wide range of operating conditions
on the breakage of agglomerates and to analyze the interplay between
breakage and other phenomena.

%=====================================================================
%=====================================================================
\section{Outline of the Thesis}
\label{sec:Outline}
\fancyhead[OL]{Outline of the Thesis}
%=====================================================================
%=====================================================================

This thesis is organized in ten chapters. The next chapter provides a
literature review to identify and analyze modeling strategies for the
structure of agglomerates, the breakage due to fluid stresses, and the
breakage by wall impacts.

 Chapters~\ref{cha:continuous_phase}~and~\ref{cha:disperse_phase}
 present a theoretical background on the description of the continuous
 and the disperse phase in the adopted Euler--Lagrange framework,
 respectively.  The proposed models for the structural features of
 agglomerates, the breakage by fluid-induced stresses, and the
 wall-impact breakage are introduced in
 Chapters~\ref{cha:agglomerate_structure},~\ref{cha:modeling_breakup_fluidstress},~and~\ref{cha:modeling_wallimpact_breakage},
 respectively. Subsequently, the numerical methods used for computing
 the continuous and the particle phase are explained in
 Chapter~\ref{cha:numerical_methodology}. In
 Chapter~\ref{cha:results_breakup_fluid}, LES predictions of
 particle-laden flows in a generic powder disperser, a square duct,
 and pipe bends relevant for dry powder inhalers are carried out and
 the results are discussed in detail mainly focusing on issues
 relevant for the breakage phenomenon.  Finally,
 Chapter~\ref{cha:conclusions} summarizes the major findings and
 conclusions from this work, and gives an outlook on related future
 research possibilities.

